For $z\geq\ell$ write $\pi_z$ for the law of $Y_{N_\ell}$ started at $z$, and for
$y\in I_i(\ell)$ write $\mu_y$ for the law of the exit state $Y_\varsigma$. The proof is an
induction on $i$, the claim being $\TV(\pi_y\|\pi_{y'})\leq(1-\omega)^{i}$ for
$y,y'\in I_i(\ell)$. At $i=0$ the descent has already ended, so $\pi_y$ and $\pi_{y'}$ are
Dirac masses and their distance is at most $1$.

At $i\geq1$, Lemma~\ref{lem:doubling_doeblin} makes $\mu_y$ a probability carried by
$I_{i-1}(\ell)$ with $\mu_y\geq\omega\varpi_i$, so
$\mu_y=\omega\varpi_i+(1-\omega)\widetilde\mu_y$ for a probability $\widetilde\mu_y$ on that
block --- the minorization is exactly what makes the remainder non-negative. The strong Markov
property at $\varsigma$, a stopping time bounded by $N_\ell$, gives
$\pi_y=\sum_{z\in I_{i-1}(\ell)}\mu_y(z)\pi_z$. Subtracting two such identities cancels the
common part $\omega\sum_z\varpi_i(z)\pi_z$ and leaves
$\pi_y-\pi_{y'}=(1-\omega)\sum_{z,z'}\widetilde\mu_y(z)\widetilde\mu_{y'}(z')(\pi_z-\pi_{z'})$,
a convex combination; the triangle inequality for the total-variation norm and the inductive
hypothesis bound it by $(1-\omega)^{i}$. That is \eqref{eq:doubling_coupling}.
